%%%%%%%%%%%%%%%%%%%%%%%%%%%%%%%%%%%%%%%%%%%%%%%%%%%%%%%%%%%%%%%%%%%%%%
% Lớp 11 | Xác suất | Biến cố giao và quy tắc nhân xác suất
% Dạng 1D9?1-2: Xác định và đếm số phần tử biến cố giao
% Số câu: 8  (sinh tự động bởi tach_cau.py; từ nay thầy sửa trực tiếp file này)
%%%%%%%%%%%%%%%%%%%%%%%%%%%%%%%%%%%%%%%%%%%%%%%%%%%%%%%%%%%%%%%%%%%%%%

%%%%%%%%%%%%%%%% TRẮC NGHIỆM 4 PHƯƠNG ÁN %%%%%%%%%%%%%%%%

% nguồn: Ok_TracNghiemLop11_Lan1
\begin{ex}%[1D9N1-2]%[ma:9452c630]
	Một hộp đựng 25 tấm thẻ cùng loại được đánh số từ 1 đến 25. Rút ngẫu nhiên một tấm thẻ trong hộp. Xét các biến số P: \lq\lq Số ghi  trên tấm thẻ là số chia hết cho 4\rq\rq \, và Q: \lq\lq Số ghi trên tấm thẻ là số chia hết cho 6\rq\rq. Nội dung của biến cố giao PQ là gì?
	\choice 
		 {$\left\{3;6;9;12;15;18;24\right\}$}
		 {$\{1;2;4\}$}
		 {$\left\{12;16;20;24\right\}$}
		 {\True $\{12;24\}$}
	 \loigiai{
	Bội chung của  $4$ và $6$ là $\left\{12;24\right\}$.}
\end{ex}

% nguồn: Ok_TracNghiemLop11_Lan1
\begin{ex}%[1D9H1-2]%[ma:78f81675]
	Gieo hai con xúc xắc cân đối và đồng chất. Gọi $A$ là biến cố \lq\lq Tổng số chấm xuất hiện trên hai con xúc xắc bằng $4$\rq\rq, $B$ là biến cố \lq\lq Tích số chấm xuất hiện trên hai con xúc xắc bằng  $4$\rq\rq. Tập hợp mô tả biến cố $AB$ là
	 \choice
		{\True $\left\{(2; 2)\right\}$} 
		{$\left\{(1; 3), (3; 1), (2; 2)\right\}$} 
		{$\left\{(1;  3), (3; 1)\right\}$}
		 {$\left\{(3; 1), (2; 2)\right\}$}
	 \loigiai{
		Tập hợp mô tả các  biến cố $A=\left\{(1; 3), (3; 1), (2; 2)\right\}$, $B=\left\{(1; 4), (4; 1), (2; 2)\right\}$.\\
		Suy ra $AB=\left\{(2;  2)\right\}$.
	} 
\end{ex}

% nguồn: Ok_TracNghiemLop11_Lan1
\begin{ex}%[1D9H1-2]%[ma:ae8a7261]
	Trên giá sách có  các quyển vở không nhãn xếp cạnh nhau với bề ngoài, khối lượng và kích thước giống hệt nhau, trong đó có $5$ quyển ghi môn Toán, $5$ quyển ghi môn Ngữ Văn và $3$ quyển ghi môn Tiếng Anh. Lấy ngẫu nhiên hai quyển vở. Xét các biến cố:\\
	$M\colon$ \lq\lq  Trong hai quyển vở được lấy, chỉ có $1$ quyển ghi môn Tiếng Anh\rq\rq;\\
	$N\colon$ \lq\lq Trong hai quyển vở được lấy, chỉ có $1$ quyển  ghi môn Ngữ Văn\rq\rq.\\
	Khi đó, biến cố giao của hai biến cố $M$  và $N$ là:
	 \choice
		{\lq\lq Hai quyển vở được lấy ghi  cùng một môn\rq\rq}
		{\lq\lq  Hai quyển vở được lấy ghi hai môn khác nhau\rq\rq}
		{ \True \lq\lq Trong hai quyển vở được lấy, một quyển  ghi môn Tiếng Anh và một quyển ghi môn Ngữ Văn\rq\rq}
		{\lq\lq Hai quyển vở  được lấy có ít nhất một quyển ghi môn Tiếng Anh\rq\rq}
	\loigiai{Biến cố giao của hai biến cố $M$ và $N$ là  \lq\lq Trong hai quyển vở được lấy, một quyển ghi môn Tiếng Anh và một quyển ghi môn Ngữ Văn\rq\rq.
	 }
\end{ex}

% nguồn: Ok_TracNghiemLop11_Lan1
\begin{ex}%[1D9H1-2]%[ma:b50ee1dd]
	Gieo một con xúc xắc cân đối và đồng chất. Gọi $A$ là biến cố \lq\lq Số chấm xuất hiện trên con xúc xắc là một số chẵn\rq\rq, $B$ là biến  cố \lq\lq Số chấm xuất hiện trên con xúc xắc chia hết cho $3$\rq\rq. Tập hợp mô tả biến cố $AB$ là
	 \choice
		{$\left\{2; 3; 4;  6\right\}$}
		{$\left\{2; 4; 6\right\}$} 
		{$\left\{3; 6\right\}$} 
		 {\True $\{6\}$}
	\loigiai{ 
		Tập hợp mô tả các biến cố  $A=\left\{2; 4; 6\right\}$, $B=\left\{3; 6\right\}$.\\
		Từ đó  suy ra $AB=\{6\}$.
	 }
\end{ex}

% nguồn: Ok_TracNghiemLop11_Lan1
\begin{ex}%[1D9H1-2]%[ma:092dc599]
	Gieo hai con xúc xắc cân đối và đồng chất. Gọi $A$ là biến cố \lq\lq Tổng số chấm xuất hiện trên hai con xúc xắc bằng $5$\rq\rq, $C$ là biến cố \lq\lq Có ít nhất một con xúc xắc xuất hiện mặt $1$ chấm\rq\rq. Tập hợp mô tả biến cố $AC$  là
	\choice 
		 {$\left\{(1; 4), (2; 3)\right\}$}
		{\True $\left\{(1; 4), (4;  1)\right\}$}
		{$\left\{(4; 1), (3;  2)\right\}$}
		{$\left\{(1; 5), (2; 3)\right\}$} 
	\loigiai{ 
		Tập hợp mô tả  các biến cố
		\[ 
		A=\left\{(1; 4), \; (4; 1),  \; (2; 3), \; (3; 2)\right\};
		\] 
		và 
		 \[
		C=\left\{(1;  6), \; (6; 1), \; (1; 5), \; (5; 1), \; (1; 4), \; (4; 1), \; (1; 3), \;(3; 1), \; (1; 2), \; (2; 1)), \; (1; 1)\right\}.
		\] 
		Từ đó suy ra $AC=\left\{(1; 4), \;  (4; 1)\right\}$.
	 }
\end{ex}

% nguồn: Ok_TracNghiemLop11_Lan1
\begin{ex}%[1D9H1-2]%[ma:76b07924]
	Xét phép thử \lq\lq Bạn thứ nhất gieo đồng xu, sau đó bạn thứ hai gieo con xúc xắc\rq\rq. Gọi $A$ là biến cố \lq\lq Đồng xu xuất hiện mặt sấp\rq\rq,  $B$ là biến cố \lq\lq Con xúc xắc xuất hiện mặt $6$ chấm''. Tập hợp mô tả biến cố $AB$ là
	 \choice
		{$\left\{(S, 4),(S, 5),(S, 6)\right\}$} 
		{$\left\{(S,  6),(N, 6)\right\}$}
		 {$\left\{(N, 6)\right\}$}
		 {\True $\left\{(S, 6)\right\}$}
	\loigiai{ 
		Tập hợp mô tả các biến cố $A=\left\{(S, 1),(S, 2),(S,  3),(S, 4),(S, 5),(S, 6)\right\}$, $B=\left\{(S, 6),(N, 6)\right\}$.\\
		Từ  đó suy ra $AB=\left\{(S, 6)\right\}$.
	 }
\end{ex}


%%%%%%%%%%%%%%%% ĐÚNG/SAI %%%%%%%%%%%%%%%%

% nguồn: Ok_dungsai_lop11_Lan1
\begin{ex}%[1D9H1-2]%[ma:c03063f8]
 Gieo hai con xúc xắc cân đối và đồng  chất
 \begin{itemize}
  \item $A$  là biến cố \lq\lq Tổng số chấm xuất hiện trên hai con xúc xắc bằng $5$\rq\rq.
  \item $B$  là biến cố \lq\lq Tích số chấm xuất hiện trên hai con xúc xắc bằng $6$\rq\rq.
  \item $C$ là biến cố \lq\lq Có ít nhất một con xúc xắc  xuất hiện mặt $1$ chấm\rq\rq.
  \item  $D$ là biến cố \lq\lq Số chấm xuất hiện trên con xúc xắc thứ nhất là $3$\rq\rq.
 \end{itemize}
 Khi đó 
  \choiceTF
  {$D=\{(3;1);(3;2);(3;3);(3;4); (3;5)\}$} 
  {\True $AD=\{(3;2)\}$} 
  {\True $BD=\{(3;2)  \}$}
  {\True  $CD=\{(3;1) \}$}
 \loigiai{ 
  \begin{itemize}
    \item Biến cố $A=\{(1;4);(4;1);(2;3);(3;2)\}$.
   \item Biến  cố $B=\{(1;6);(6;1);(2;3);(3;2)\}$.
   \item Biến cố $C=\{(1;6);(6;1);(1;5);(5;1);(1;4);(4;1);(1;3);(3;1);(1;2);(2;1);(1;1)\}$. 
   \item Biến cố $D=\{(3;1);(3;2);(3;3);(3;4); (3;5);  (3;6)\}$.
  \end{itemize}
  Khi đó, 
  \begin{itemchoice}
   \itemch Biến cố $D=\{(3;1);(3;2);(3;3);(3;4); (3;5);  (3;6)\}$.
    \itemch Biến cố $AD=\{(3;2)\}$.
   \itemch Biến  cố $BD=\{(3;2) \}$.
   \itemch Biến  cố $CD=\{(3;1) \}$.
  \end{itemchoice}
 } 
\end{ex}

% nguồn: Ok_dungsai_lop11_Lan1
\begin{ex}%[1D9H1-2]%[ma:cb5f1bee]
 Gieo hai con xúc  xắc cân đối và đồng chất.
 \begin{itemize}
   \item $A$ là biến cố \lq\lq Tổng số chấm xuất hiện trên hai con xúc xắc bằng $5$\rq\rq.
  \item $B$ là biến cố  \lq\lq Tích số chấm xuất hiện trên hai con xúc xắc bằng $6$\rq\rq.
  \item $C$ là biến cố \lq\lq Có ít nhất một con xúc xắc xuất hiện  mặt $1$ chấm\rq\rq.
 \end{itemize}
  Khi đó
 \choiceTF 
  {\True  $A=\{(1;4);(4;1);(2;3);(3;2)\}$}
  {\True  $AC=\{(1;4);(4;1)\}$}
   {$AB=\{(2;3)\}$}
  {\True  $BC=\{(1;6);(6;1)\}$}
  \loigiai{
  \begin{itemize}
   \item Biến  cố $A=\{(1;4);(4;1);(2;3);(3;2)\}$.
    \item Biến cố $B=\{(1;6);(6;1);(2;3);(3;2)\}$.
   \item Biến cố $C=\{(1;6);(6;1);(1;5);(5;1);(1;4);(4;1);(1;3);(3;1);(1;2);(2;1);(1;1)\}$. 
  \end{itemize}
  Khi  đó,
  \begin{itemchoice}
   \itemch Biến cố  $A=\{(1;4);(4;1);(2;3);(3;2)\}$.
   \itemch Biến cố $AC=\{(1;4);(4;1)\}$. 
   \itemch Biến cố $AB=\{(2;3);(3;2)\}$. 
    \itemch Biến cố $BC=\{(1;6);(6;1)\}$.
  \end{itemchoice}
  }
\end{ex}
